\chapter{Methodology}
\label{ch:methodology}

This chapter describes the proposed system in detail. The overall architecture is presented first (Section~\ref{sec:overview}), followed by the tiered classifier and its two backbones (Section~\ref{sec:tiered}), the confidence-based routing logic (Section~\ref{sec:routing}), the uncertainty-quantification methods (Section~\ref{sec:uncertainty}), the calibration and conformal-prediction components (Sections~\ref{sec:calibration}--\ref{sec:conformal}), the synthetic-augmentation pipeline (Section~\ref{sec:synthetic}), and finally the explainability and software-engineering aspects (Sections~\ref{sec:explain}--\ref{sec:software}).

\section{System Overview}
\label{sec:overview}
The proposed tiered chest-radiograph classification architecture is shown in Figure~\ref{fig:architecture}. The system is organised around tiered routing logic that activates a lightweight or a deep model according to the estimated difficulty of each image. An input radiograph is first preprocessed and passed through the Tier-1 screening model; if Tier~1 is sufficiently confident the decision is returned immediately, and otherwise the image is escalated to the heavier Tier-2 model, which is executed under Monte Carlo Dropout and test-time augmentation and wrapped in a conformal predictor. The final output is not merely a label but a label together with a confidence, an uncertainty estimate, and a conformal prediction set.

% System architecture diagram (TikZ). Styles defined in main.tex preamble.
\begin{figure}[htbp]
\centering
\resizebox{\textwidth}{!}{%
\begin{tikzpicture}
  \node[io] (input) at (0,0) {Chest\\radiograph};
  \node[proc, right=0.7cm of input] (pre) {Preprocessing\\[1pt]{\scriptsize CLAHE, $224{\times}224$,}\\[-1pt]{\scriptsize ImageNet norm.}};
  \node[light, right=0.7cm of pre] (t1) {\textbf{Tier 1}\\MobileNetV2\\[1pt]{\scriptsize $\rightarrow (p_1,\ \text{conf.})$}};
  \node[decision, right=0.9cm of t1] (dec) {conf.\\$\geq \tau$?};
  \node[io, right=3cm of dec] (out) {\textbf{Output}\\[1pt]{\scriptsize prediction,}\\[-1pt]{\scriptsize uncertainty,}\\[-1pt]{\scriptsize conformal set}};
  \node[heavyblk, below=1.7cm of dec] (t2) {\textbf{Tier 2}\\EfficientNet-B4 / Ark+ Swin\\[1pt]{\scriptsize MC-Dropout ($T{=}20$), TTA ($\times 10$)}};
  \node[proc] (cp) at (out |- t2) {Conformal\\prediction\\[1pt]{\scriptsize set, cov.\ $0.95$}};

  \draw[flow] (input) -- (pre);
  \draw[flow] (pre) -- (t1);
  \draw[flow] (t1) -- (dec);
  \draw[flow] (dec) -- node[flowlab, above, pos=0.4]{accept ($\approx$74\%)} (out);
  \draw[flow] (dec) -- node[flowlab, right]{escalate ($\approx$26\%)} (t2);
  \draw[flow] (t2) -- (cp);
  \draw[flow] (cp) -- (out);
\end{tikzpicture}%
}
\caption{Overview of the tiered classification and confidence-based routing system. A lightweight Tier~1 screening model (MobileNetV2) resolves confident cases directly; only low-confidence images are escalated to the heavier Tier~2 backbone, which is run with Monte Carlo Dropout and test-time augmentation and wrapped in a conformal predictor.}
\label{fig:architecture}
\end{figure}


\section{Tiered Classification System}
\label{sec:tiered}
The first stage of the system (Tier~1) is the computationally efficient MobileNetV2 model \cite{sandler2018mobilenetv2}. For an input image $x$, Tier~1 produces a class-probability vector $p_1(x) = \mathrm{softmax}(f_1(x))$, where $f_1$ denotes the Tier-1 network. The model's confidence is taken to be the probability of its predicted class,
\begin{equation}
    c_1(x) = \max\big(p_1(x)_0,\ p_1(x)_1\big),
\end{equation}
which lies in $[0.5, 1]$ for the binary pneumothorax-versus-no-finding task. When $c_1(x)$ exceeds the current escalation threshold $\tau$, the final prediction is produced by Tier~1 alone.

If Tier~1 is not sufficiently confident ($c_1(x) < \tau$), the image is escalated to the second stage (Tier~2), which has stronger feature-representation capabilities. Two alternative Tier-2 backbones are studied: EfficientNet-B4 \cite{tan2019efficientnet}, a compound-scaled convolutional network, and the Ark+ foundation model \cite{ma2023foundation} with a Swin Transformer backbone \cite{liu2021swin}. Because Tier~1 resolves the large majority of cases, the heavy Tier-2 model is invoked only on the difficult minority, which is the source of the system's efficiency.

\section{Confidence-Based Routing}
\label{sec:routing}
The routing decision is governed by the escalation threshold $\tau$. Two regimes are considered. In the \emph{static} regime, $\tau$ is fixed at a value chosen on the validation set. In the \emph{dynamic} regime, $\tau$ adapts online to the stream of recent confidences: the router maintains a sliding window of the last $W$ Tier-1 confidence values and nudges $\tau$ up or down by a small step $\delta$ so as to keep the escalation rate stable in the face of distribution drift. The complete routing logic, including the escalation decision, the high-variance review flag, and the conformal abstention rule, is shown in Figure~\ref{fig:routing}.

% Confidence-based routing flowchart (TikZ). Styles defined in main.tex preamble.
\begin{figure}[htbp]
\centering
\begin{tikzpicture}[node distance=0.95cm]
  \node[light] (t1) {Tier-1 inference $\rightarrow$ confidence $c$};
  \node[proc, below=of t1] (win) {Update sliding window (last 50);\\adapt threshold $\tau$ around $0.75$ by $\delta=0.05$};
  \node[decision, below=of win] (d1) {$c \geq \tau$?};
  \node[light, right=2.6cm of d1] (accept) {Accept Tier-1\\prediction};
  \node[heavyblk, below=1.15cm of d1] (t2) {Tier-2 inference\\{\scriptsize MC-Dropout ($T{=}20$) + TTA ($\times 10$) $\rightarrow$ mean prob, variance $\sigma^2$}};
  \node[decision, below=of t2] (d2) {$\sigma^2 > 0.08$?};
  \node[proc, right=2.6cm of d2] (flag) {Flag for\\radiologist review};
  \node[proc, below=of d2] (cp) {Conformal set: if $|\,C(x)\,| > 1$, abstain (return both labels)};
  \node[io, below=of cp] (end) {Final decision $+$ confidence $+$ uncertainty $+$ conformal set};

  \draw[flow] (t1) -- (win);
  \draw[flow] (win) -- (d1);
  \draw[flow] (d1) -- node[flowlab, above]{yes} (accept);
  \draw[flow] (d1) -- node[flowlab, right]{no} (t2);
  \draw[flow] (t2) -- (d2);
  \draw[flow] (d2) -- node[flowlab, above]{yes} (flag);
  \draw[flow] (d2) -- node[flowlab, right]{no} (cp);
  \draw[flow] (cp) -- (end);
  % Merge the two "yes" branches into the final node along right-hand rails that
  % clear the right-column boxes (the accept line must not pass through the flag box).
  \draw[flow] (accept.east) -- ++(1.6,0) |- (end.east);
  \draw[flow] (flag.east) -- ++(0.7,0) |- (end.east);
\end{tikzpicture}
\caption{Confidence-based routing and escalation logic. The escalation threshold is adapted online from a sliding window of recent confidences; escalated cases are resolved by the Tier-2 backbone under Monte Carlo Dropout and test-time augmentation, and high-variance predictions are flagged for human review.}
\label{fig:routing}
\end{figure}


Formally, let $\bar{c}_t$ be the mean confidence over the current window at time $t$. The dynamic threshold update is
\begin{equation}
    \tau_{t+1} =
    \begin{cases}
        \min(\tau_{\max},\ \tau_t + \delta) & \text{if } \bar{c}_t \text{ is high (few hard cases),}\\
        \max(\tau_{\min},\ \tau_t - \delta) & \text{otherwise,}
    \end{cases}
\end{equation}
with $\tau$ initialised at the base value and clamped to a sensible range. In the experiments the base threshold is $0.75$, the window size is $W=50$, and the step is $\delta=0.05$ (Chapter~\ref{ch:experimental_setup}). This adaptive rule keeps the proportion of escalated images near a target operating point, which in our experiments settles at roughly one quarter of the test set.

\section{Uncertainty Quantification}
\label{sec:uncertainty}
Uncertainty is estimated with two complementary methods that are applied to the Tier-2 model.

\paragraph{Monte Carlo Dropout.} Following Gal and Ghahramani \cite{gal2016dropout}, the dropout layers of the Tier-2 network are kept active at inference time and $T$ stochastic forward passes are performed for the same input. Let $p^{(t)}(x)$ be the probability vector from the $t$-th pass. The predictive mean and the per-class variance are
\begin{equation}
    \bar{p}(x) = \frac{1}{T}\sum_{t=1}^{T} p^{(t)}(x),
    \qquad
    \sigma^2(x) = \frac{1}{T}\sum_{t=1}^{T}\big(p^{(t)}(x) - \bar{p}(x)\big)^2 .
\end{equation}
The variance $\sigma^2(x)$ of the positive class is used as the epistemic-uncertainty signal; predictions whose variance exceeds a flagging threshold are marked for human review. Importantly, only the dropout layers are stochastic --- the batch-normalisation statistics are kept in evaluation mode --- so that the variance reflects model uncertainty rather than fluctuating normalisation.

\paragraph{Test-Time Augmentation.} In addition, the input is subjected to a set of label-preserving augmentations (small rotations, translations, and brightness changes), and the predictions over the augmented copies are averaged. The dispersion of these predictions measures the sensitivity of the decision to plausible perturbations of the input \cite{wang2019aleatoric}. Monte Carlo Dropout and test-time augmentation are combined by tiling the $T$ stochastic passes across the augmented views in a single batched forward computation, which keeps the additional cost tractable on a GPU.

\section{Calibration}
\label{sec:calibration}
Because the raw softmax outputs of deep networks are typically over-confident \cite{guo2017calibration}, a post-hoc temperature-scaling step is applied. A single scalar temperature $T_{\text{scale}}$ is fitted on the validation split by minimising the negative log-likelihood, and the logits $z$ are rescaled as $\mathrm{softmax}(z / T_{\text{scale}})$ at inference time. Calibration quality is summarised by the Expected Calibration Error (ECE), defined in Chapter~\ref{ch:experimental_setup}.

\section{Conformal Prediction Integration}
\label{sec:conformal}
To equip the predictions with a finite-sample coverage guarantee, the Conformal Prediction methodology \cite{vovk2005algorithmic, shafer2008tutorial} is applied. Using a dedicated calibration split that is disjoint from both the training and the test data, a non-conformity score is computed for each calibration example as
\begin{equation}
    s_i = 1 - \hat{f}(x_i)_{y_i},
\end{equation}
where $\hat{f}(x_i)_{y_i}$ is the calibrated model probability assigned to the true class $y_i$. For a chosen error tolerance $\alpha$, the threshold $\hat{q}$ is taken to be the $\lceil (n+1)(1-\alpha) \rceil / n$ empirical quantile of the calibration scores $\{s_i\}$. At test time, the prediction set for a new image $x_{\text{test}}$ is
\begin{equation}
    C(x_{\text{test}}) = \{\, y : 1 - \hat{f}(x_{\text{test}})_y \le \hat{q} \,\}.
\end{equation}
Under exchangeability, the probability that this set contains the true label is at least $1-\alpha$. The size of $C(x_{\text{test}})$ is itself informative: a singleton set is a confident decision, whereas a set containing both labels is an explicit, calibrated abstention. The target coverage used throughout is $1-\alpha = 0.95$.

\section{Synthetic Data Augmentation}
\label{sec:synthetic}
The pneumothorax class is a minority in the training data (Chapter~\ref{ch:experimental_setup}). To enrich it, a Stable-Diffusion synthetic-augmentation pipeline is used, shown in Figure~\ref{fig:sd_pipeline}. A latent-diffusion model \cite{rombach2022high} generates additional pneumothorax-positive radiographs, and each synthetic image is admitted to the Tier-2 training set only if its Fr\'echet Inception Distance to the real distribution lies below a quality threshold. The accepted synthetics oversample the minority class without the mode collapse and instability associated with adversarial generators. The contribution of this pipeline is isolated by two ablations: one that removes the synthetic data (A8) and one that removes all augmentation (A9).

% Stable-Diffusion synthetic augmentation pipeline (TikZ). Styles in main.tex preamble.
\begin{figure}[htbp]
\centering
\resizebox{\textwidth}{!}{%
\begin{tikzpicture}
  \node[io] (real) at (0,0) {Real minority-class\\images\\[1pt]{\scriptsize (pneumothorax-positive,}\\[-1pt]{\scriptsize $\approx$16.7\% prevalence)}};
  \node[heavyblk, right=0.9cm of real] (sd) {Stable Diffusion\\{\scriptsize latent diffusion,}\\[-1pt]{\scriptsize 50 steps, guidance 7.5}};
  \node[decision, right=0.9cm of sd] (fid) {FID\\$\leq 50$?};
  \node[proc, above=1.0cm of fid] (reject) {Reject};
  \node[light, right=1.4cm of fid] (aug) {Augmented\\training set\\[1pt]{\scriptsize real $+$ accepted}\\[-1pt]{\scriptsize synthetic images}};
  \node[heavyblk, right=0.9cm of aug] (train) {Tier-2\\training};

  \draw[flow] (real) -- (sd);
  \draw[flow] (sd) -- (fid);
  \draw[flow] (fid) -- node[flowlab, right]{no} (reject);
  \draw[flow] (fid) -- node[flowlab, above]{yes} (aug);
  \draw[flow] (aug) -- (train);
\end{tikzpicture}%
}
\caption{Stable-Diffusion synthetic augmentation pipeline for the minority (pneumothorax) class. Generated images are admitted to the Tier-2 training set only when their Fr\'echet Inception Distance to the real distribution is below the quality threshold. Ablations A8 (no synthetic data) and A9 (no augmentation at all) quantify the contribution of this pipeline.}
\label{fig:sd_pipeline}
\end{figure}


\section{Explainability}
\label{sec:explain}
For each prediction, a HiResCAM \cite{draelos2020hirescam} saliency map is produced for the model that made the decision (Tier~1 or Tier~2). HiResCAM was chosen over the more common Grad-CAM \cite{selvaraju2017gradcam} because it is provably faithful to the model's computation for the architectures used here and avoids the diffuse, sometimes misleading activations that Grad-CAM can produce. The saliency map is computed on the un-normalised image so that the highlighted regions correspond to anatomically meaningful structures rather than to artefacts of the input normalisation.

\section{Software Architecture}
\label{sec:software}
The system is implemented as a layered application that separates the domain logic (models, routing, uncertainty, evaluation) from the application services and the delivery infrastructure (a FastAPI back end and a React front end). Established design patterns --- a factory for model construction, a strategy for the routing policy, an observer for training callbacks, and a repository for data access --- keep the components decoupled, and the layer boundaries are enforced automatically by an import-contract check in continuous integration. This engineering discipline is what makes the fifteen-configuration ablation study of Chapter~\ref{ch:results} reproducible: every configuration is expressed as data and evaluated by a single unified evaluator.
